\section{Arutelu} \label{arutelu}

Käesolev peatükk arutleb eksperimendi tulemuste üle laiemas kontekstis, käsitleb TI kasutamise praktilisi aspekte, piiranguid ja annab soovitusi.

\subsection{TI tugevused lõputöö kirjutamisel} \label{subsec:tugevused}

Eksperimendi tulemused kinnitavad, et kaasaegsed TI mudelid koos agentsete tööriistadega suudavad pakkuda märkimisväärset abi lõputöö kirjutamise protsessis. Käesoleva uurimuse katsetes demonstreeris Claude Opus 4.6 + Coworki kombinatsioon eriti veenvalt, et tipptasemel mudel koos agentse tööriistaga suudab genereerida tervikliku 61-leheküljelise eestikeelse lõputöö ka ühe kasutajaviibaga ilma iteratiivse suunamiseta. Peatükis~\ref{subsec:opus47-juhtum} esitatud iteratiivne juhtumiuuring uuema mudeliga Claude Opus 4.7 + Claude Code näitas omakorda, et kui kasutaja juhendamine on aktiivne ja seanss kestab ühe tööpäeva, tõuseb saavutatav kvaliteet märgatavalt. Valminud 57-leheküljeline töö sisaldas tegelikku, kompileeritavat Pythoni ja TypeScripti lähtekoodi, 32 kontrollitavat viidet ja alla 0,1 kirjaviga lehekülje kohta. Peamised tugevused jagunevad järgmiselt.

\textbf{Struktureerimine ja kavandamine.} TI mudelid on eriti tugevad teksti struktureerimisel. Kõik neli testitud kombinatsiooni (arvestades GPT-5.4 Codex-laienduse kaudu saadud tulemust) suutsid luua loogilise ja akadeemilistele nõuetele vastava peatükkide jaotuse ainuüksi ühe kasutajaviibaga. Eriti kasulik on TI kasutamine töö algfaasis, kui üliõpilane vajab abi sisukorra ja peatükkide ülesehituse kavandamisel. Mudelid suudavad vaid minimaalse juhise põhjal luua detailse sisukorra koos soovitusliku leheküljemahuga iga peatüki jaoks.

\textbf{Mustandi loomine.} TI suudab genereerida tekstimustandeid, mis on piisavalt hea kvaliteediga, et neid edasi töödelda. See kiirendab kirjutamisprotsessi oluliselt, kuna üliõpilane ei pea alustama tühjalt lehelt, vaid saab toimetada ja täiendada olemasolevat teksti. Esialgse kasutuskogemuse põhjal võib TI mustandi kasutamine säästa hinnanguliselt 50--80\% kirjutamisele kuluvast ajast, kuigi täpse ajavõidu mõõtmine vajab edaspidi süstemaatilisi kasutajakatseid.

\textbf{Keeleline toimetamine.} Kõik neli mudelit suudavad keelelist toimetamist toetada -- parandada grammatikavigu, lauseehitust ja ühtlustada stiili, kuid täpsus erineb mudeliti. See on eriti väärtuslik eesti keele kontekstis, kus spetsiifilisi akadeemilise toimetamise teenuseid on vähe.

\textbf{\LaTeX{}-vormindamine.} Oluline praktiline eelis on mudelite võimekus genereerida teksti otse \LaTeX{}-vormingus, mis hõlmab korrektset viitamissüntaksit, tabeleid, jooniste viiteid ja matemaatilisi valemeid. See säästab üliõpilasele märkimisväärset vormistusaega. Võrdluseks, Microsoft Wordi puhul peab üliõpilane vormistuse -- stiilid, ristviited, pealkirjade numeratsioon ja bibliograafia -- enamasti käsitsi seadistama, kuna TI mudelid ei saa \texttt{.docx} faili sisemist struktuuri samamoodi manipuleerida nagu tekstipõhist \LaTeX{}-lähtekoodi. See teeb Wordi TI-toetatud akadeemilise kirjutamise kontekstis oluliselt aeganõudvamaks ja veaohtlikumaks valikuks.

\subsection{TI puudujäägid ja piirangud} \label{subsec:puudujaagid}

Vaatamata edusammudele on TI mudelitel lõputöö kirjutamise kontekstis endiselt olulised puudujäägid. Tulemuste tõlgendamisel tuleb eriti arvestada kolme metoodilist piirangut: eksperimendis võrreldi \emph{mudeli ja agentse tööriista kombinatsioone}, mitte mudeleid isoleeritult; iga kombinatsiooni kohta tehti vaid üks jooks; ning eksplitsiitset kontrollgruppi -- ei inimkirjutatud lõputöö võrdluspunkti ega mudelite mitteagentset baasväljundit (peale GPT-5.4 vestlusliidese erandi) -- ei kasutatud. Seetõttu näitavad tulemused konkreetsete töövoogude käitumist selles katses, mitte lõplikku mudelite üldist paremusjärjestust ega kvaliteeti absoluutse bakalaureusetöö läve suhtes. Süstemaatiline hinnang nõuaks iga kombinatsiooni kohta mitut kordusjooksu, suuremat valimit ja inimkirjutatud lõputöö kontrollvõrdlust.

\textbf{Viidete usaldusväärsus.} Üks tõsisemaid probleeme on viidete hallutsineerimine, kuid katsetulemused näitavad, et hallutsineerimismäärad varieeruvad mudelite vahel oluliselt. Käesolevas eksperimendis kontrollis töö autor iga viite käsitsi, et veenduda, kas allikas tegelikult eksisteerib ning kas pealkiri, autorid, väljaandmise aasta ja URL/DOI vastavad tegelikkusele. Tulemused olid järgmised: GPT-5.4 + Codex 0\% (kõik 22 viidet täielikult korrektsed), Opus 4.6 + Cowork 6\% (35-st 2 viidet väikese metaandmete veaga), Gemini 3.1 Pro + Antigravity 11\% (9-st 1 viide vale autoriga) ning Sonnet 4.6 + Cowork 22\% (23-st 5 viidet välja mõeldud autoriloendiga). Sonnet 4.6 hallutsinatsioonimuster oli tähelepanuväärne -- artiklite pealkirjad ja arXiv-identifikaatorid olid valdavalt korrektsed, kuid mudel ``täitis'' autoriloendeid välja mõeldud nimedega või kombineeris kahe eri paberi metaandmed (näiteks viide ``JITBot'' Pornprasit ja Tantithamthavorn 2022, kus tegelikult oli paber ``JITLine'' MSR 2021 või JITBot ASE 2020 hoopis teiste autoritega). Need tulemused on kooskõlas varasemate uurimustega: Chelli jt \cite{chelli2024hallucination} tuvastasid GPT-4 hallutsinatsioonimääraks 28,6\%; Mugaanyi jt \cite{mugaanyi2024citation} leidsid, et DOI-de täpsus on valdkonniti väga erinev (loodusteadustes 32,7\% vs humanitaarteadustes 8,5\%), mis viitab sellele, et arvutiteaduse valdkonnas võivad tulemused olla mõnevõrra paremad. Kim jt \cite{kim2025geographic} on näidanud, et hallutsinatsioonid on eriti levinud väiksemate keelte kontekstis. Iga TI-genereeritud viidet peab seega kontrollima inimene -- ka madala üldise hallutsinatsioonimääraga mudelite puhul, sest pealkirja olemasolu ei tähenda, et autorid või muud metaandmed oleksid samuti korrektsed.

\textbf{Sisuline sügavus.} TI-genereeritud tekstid kipuvad olema pealiskaudsed. Mudelid kordavad üldtuntud seisukohti, kuid ei suuda luua originaalset analüüsi ega teha uudseid järeldusi. See on eriti problemaatiline tulemuste arutelu puhul, kus eeldatakse, et autor demonstreerib kriitilist mõtlemist ja seob tulemused laiema kontekstiga kokku.

\textbf{Eesti keele tehniline terminoloogia.} Kuigi mudelite eesti keele oskus on oluliselt paranenud, esineb endiselt probleeme tehnilise terminoloogiaga. Mudelid kasutavad kohati ebaühtlast terminoloogiat (vahetavad eesti- ja ingliskeelsete terminite vahel) ning mõnikord loovad valesid termineid. Näiteks kasutati sõna ``treenima'' asemel ``koolitama'', mis ei vasta arvutiteaduse valdkonnas juurdunud eestikeelsele terminoloogiale. Barbu jt \cite{barbu2025estonian} on samuti täheldanud, et keelemudelid vajavad eesti keele spetsiifika jaoks täiendavat kohandamist, mis viitab laiemale probleemile väikeste keelte tehnilistel kasutusjuhtudel.

\textbf{Kontekstiakna piirangud.} Kuigi kontekstiaknad on oluliselt suurenenud, on need endiselt piiratud. 25-leheküljelise bakalaureusetöö tekst (ca 80\,000--100\,000 tähemärki ehk ca 40\,000--50\,000 tokenit) mahub küll kontekstiaknasse, kuid kui lisada juhendid, näidistööd ja varasemalt genereeritud osad, on kontekstiakna efektiivne kasutamine endiselt kriitiline.

\textbf{Juhiste järgimise ebaühtlus ning liidese ja agentsete tööriistade mõju.} Peatükis~\ref{eksperimendid} kirjeldatud tulemused näitavad, et sama mudel võib eri liidestes käituda sisuliselt erinevalt. GPT-5.4 vestlusliidese keeldumine ja Codex-laienduse kaudu õnnestunud failipõhine genereerimine viitavad sellele, et väljundit ei määra ainult alusmudel, vaid ka süsteemijuhised, ohutuspoliitika ja töövoo ülesehitus. Bianchi jt \cite{bianchi2023safety} kirjeldatud ``liialdatud ohutuskäitumine'' aitab seda nähtust tõlgendada: liides võib uurimusliku ülesande tõlgendada akadeemilise väärkasutuse riskina. Agentne tööriist ei tähenda seejuures ohutuspiirangutest möödumist, vaid teistsugust ülesande konteksti, kus töö jaotub väiksemateks faili- ja kompileerimispõhisteks sammudeks.

Sama oluline on eristada vormilist veenvust sisulisest usaldusväärsusest. Gemini 3.1 Pro + Antigravity töö näitas kõige selgemalt, et korrektselt vormistatud ja statistiliste tabelitega tekst võib sisaldada täielikult väljamõeldud empiirikat. Opus 4.6 + Cowork andis selles katses tugevama ja tehniliselt sügavama väljundi, kuid ka seal tuli kontrollimata testitulemusi ja viiteid käsitleda riskina. Seega ei piisa agentse tööriista edukast \LaTeX{}-vormistusest -- sisuline vastutus, andmete olemasolu ja viidete kontrollimine peavad jääma inimesele.

\subsection{Hinnakujundus ja majanduslik teostatavus} \label{subsec:hinnakujundus}

Antud peatüki kuluarvestus piirdub peatükis~\ref{subsec:opus47-juhtum} kirjeldatud iteratiivse juhtumiuuringuga (Claude Opus 4.7 + Claude Code), kuna ainult selle töövoo tokenikasutus telemetreeriti reaalajas ning seda saab esitada mõõdetud andmetena. Põhieksperimendi nelja ühe viiba kombinatsiooni (Gemini 3.1 Pro + Antigravity, GPT-5.4 + Codex, Claude Sonnet 4.6 + Cowork, Claude Opus 4.6 + Cowork) tokenitarvet eraldi ei mõõdetud, mistõttu nendele konkreetseid kuluhinnanguid ei esitata; ametlikud baashinnakirjad on esitatud peatükis~\ref{subsec:koondvordlus} (tabel~\ref{tab:mudelite-vordlus-hind}).

\begin{table}[ht]
\centering
\caption{Iteratiivse juhtumiuuringu (Opus 4.7 + Claude Code) hüpoteetiline pay-as-you-go API-kulu komponentide kaupa, kui kuutellimuse asemel oleks kasutatud ametlikku API-d (EUR, koos 24\% käibemaksuga)}
\label{tab:hinnavordlus}
\begin{tblr}{
    width=\linewidth,
    colspec={X[3,l] X[1,c]},
    rows={font=\small},
    colsep=3pt,
    row{1}={font=\bfseries\small},
    hlines,
    vlines,
    row{even}={bg=colorCellHighlight},
}
Komponent & API-kulu \\
Claude Code (Opus 4.7) -- põhigenereerimine & ca 322~€ \\
Codex (GPT-5.5) -- keelekorrektuur ja sõnastuse ühtlustamine & ca 46~€ \\
Hindamise 80 API-päringut & ca 4~€ \\
Veebiotsingu 22 kutset & alla 0,30~€ \\
\textbf{Kokku} & \textbf{ca 371~€} \\
\end{tblr}
\end{table}

Hinnad on esitatud eurodes, lähtudes mudelite pakkujate 2026. aasta mai ametlikest USD-hinnakirjadest \cite{openai2026pricing, anthropic2026pricing, google2026gemini_pricing} ja teisendatud kursi 1 EUR = 1,17 USD järgi (vt peatükk~\ref{subsec:koondvordlus}). Esitatud summad sisaldavad Eesti standardset 24\% käibemaksu, mis vastab eraisikust üliõpilase tegelikule kulule.

Iteratiivse juhtumiuuringu tegelik mõõdetud tokenikasutus jagunes järgmiselt. Claude Code (Opus 4.7) tööriistas registreeriti seansi vältel ligikaudu 5,7 tuhat vahemälust välja jäänud sisendtokenit, 8,15 miljonit vahemälu loomise tokenit, 403,0 miljonit vahemälust loetud sisendtokenit ja 2,03 miljonit väljundtokenit. Anthropicu Opus 4.7 standardhinnakirjal (sisend 5~USD/MTok, vahemälu kirjutamine 6,25~USD/MTok, vahemälust lugemine 0,50~USD/MTok, väljund 25~USD/MTok) oleks see vastanud ligikaudu 303~USD-le ehk käibemaksuga ca 322~€-le. Paralleelselt kasutati keelekorrektuuri ja sõnastuse ühtlustamise jaoks OpenAI ChatGPT Codex laienduse kaudu GPT-5.5 mudelit kogumahus ligikaudu 38,7~miljonit sisendtokenit (millest 35,7~M oli vahemälust loetud) ja 0,33~M väljundtokenit, mille hüpoteetiline API-kulu OpenAI GPT-5.5 lühikese konteksti hinnaga (sisend 5~USD/MTok, vahemälust loetud sisend 0,50~USD/MTok, väljund 30~USD/MTok) oleks olnud ca 43~USD ehk käibemaksuga ca 46~€. Lisaks tehti peatükis~\ref{subsec:opus47-juhtum} kirjeldatud 80 hindamis-API-päringut (kogukulu 3,88~USD ehk ca 4~€ käibemaksuga) ning 22 veebiotsingu kutset (OpenAI hinnakirjas 10~USD/1000 kutset ehk ca 0,22~USD ehk alla 0,30~€). Komponendid kokku liites oleks kogu iteratiivse juhtumiuuringu hüpoteetiline API-kulu olnud ligikaudu 350~USD ehk Eestis ca 371~€ käibemaksuga. Suur osa sellest summast tuleneb vahemälu korduvlugemisest agentse failipõhise töövoo vältel, mistõttu sama kasutuse katmine kuutellimuse raames jääb üliõpilase jaoks sadu eurosid soodsamaks.

Kuutellimused muudavad kulupildi oluliselt teistsuguseks: ChatGPT Plus ja Claude Pro maksavad Eesti regioonis ca 22--23~€ kuus (koos käibemaksuga), Google AI Pro 21,99~€ kuus (vt tabel~\ref{tab:mudelite-tellimused}). Praktikas saab üliõpilane teha põhilise töö 1--2 kuu tellimuse hinnaga, kui mahupiiranguid ei ületata (eeldusel, et ei soovita väga pikki sessioone järjest teha). Ka iteratiivne juhtumiuuring viidi käesolevas töös läbi olemasoleva Claude'i ja OpenAI kuutellimuse raames, mistõttu tegelik täiendav kulu üliõpilase jaoks oli marginaalne -- eelmises lõigus arvutatud ca 371~€ kirjeldab seega pigem võrdluseks pay-as-you-go hinnatabelist tuletatud teoreetilist maksimumi kui üliõpilase tegelikku kulu.

Kokkuvõttes näitab mõõdetud iteratiivse juhtumiuuringu API-kulu (ca 371~€), et isegi tipptasemel mudeli ja agentse tööriista pikem juhitud seanss jääb üldjuhul oluliselt soodsamaks kui professionaalse toimetaja teenus, mis võib bakalaureusetöö puhul maksta mitusada eurot. Kuutellimuse raames kasutamisel langeb tegelik täiendav kulu üliõpilase jaoks veelgi madalamaks ning piirdub jooksva tellimuse maksega.

\subsection{Eesti keele korrektuur ja toimetamine} \label{subsec:korrektuur}

Eraldi väärib tähelepanu TI mudelite kasutamine eestikeelse teksti korrektuuri ja toimetamise tööriistana. See on üks valdkondi, kus TI pakub kõige suuremat praktilist väärtust.

Järgmised hinnangud põhinevad autori mitteformaalsetel katsetel, mis jäid põhieksperimendi raamistikust välja ning mille täpsem metoodika vajab edaspidi süstemaatilist vormistamist. Need on esitatud suunavate tähelepanekutena.

\textbf{Grammatikakontroll.} Kõik neli mudelit suudavad tuvastada ja parandada enamiku grammatikavigadest eesti keeles. Claude Opus 4.6 oli kõige täpsem, tuvastades enamiku sisestatud vigadest; Sonnet 4.6 ja GPT-5.4 tuvastasid palju, kuid jätsid osa keerulisemaid vigu märkamata; Gemini 3.1 Pro tulemus oli nõrgim, jättes tuvastamata ka lihtsamaid morfoloogiavigu.

\textbf{Stiilikontroll.} Mudelid suudavad hinnata teksti vastavust akadeemilisele stiilile ja pakkuda parandusettepanekuid. See hõlmab liiga vaba keele tuvastamist, passiivi kasutamise soovitusi ja liigsete korduste märkamist. Claude Opus 4.6 oli selles osas kõige põhjalikum, pakkudes detailsemaid selgitusi paranduste kohta.

\textbf{Terminoloogiakontroll.} Mudelid suudavad kontrollida tehnilise terminoloogia ühtlust ja korrektsust, kuid siin on täpsus oluliselt madalam kui grammatikakontrollis. Terminoloogilist kontrolli ei saa täielikult TI-le delegeerida ja valdkonnaspetsiifiline inimekspertiis on endiselt vajalik.

\subsection{Soovitused TI tööriista valikuks} \label{subsec:soovitused}

Eksperimendi ja iteratiivse juhtumiuuringu tulemuste ning praktiliste kogemuste põhjal saab anda järgmised soovitused.

\textbf{Iteratiivse koostöö parim tulemus: Claude Opus 4.7 koos Claude Code'iga.} Peatükis~\ref{subsec:opus47-juhtum} kirjeldatud juhtumiuuring näitas, et kui kasutaja võtab aktiivse rolli ja seanss kestab ühe tööpäeva (ligikaudu kaheksa tundi), suudab Opus 4.7 + Claude Code'i kombinatsioon genereerida 57-leheküljelise lõputöö koos tegeliku kompileeritava Pythoni ja TypeScripti lähtekoodiga. Tulemus on ortograafilise täpsuse, viidete usaldusväärsuse ja sisulise terviklikkuse poolest mõõdetavalt parem kui üheski neljast ühe viiba kombinatsioonist.

\textbf{Ühe viibaga tugevaim üldtulemus: Claude Opus 4.6 koos Coworkiga.} Käesolev eksperiment näitas, et ühe kasutajaviibaga genereerimisel eristus Claude Opus 4.6 + Coworki kombinatsioon teistest olulisel määral terviklikkuse, akadeemilise stiili, tehnilise detailsuse ja sisulise sügavuse poolest, kuigi väikseim kirjavigade arv oli GPT-5.4 + Codexi väljundis. Peamine puudus on märgatavalt kõrgem hind. Kuna iga kombinatsiooni kohta tehti vaid üks jooks ja keelemudelid on mitteterministlikud, tuleb seda soovitust võtta eksploratiivse viitena, mitte lõpliku üldhinnangu kinnitusena.

\textbf{Parim hinna-kvaliteedi suhe: Claude Sonnet 4.6 koos Coworkiga.} Claude Sonnet 4.6 + Coworki kombinatsioon on kõige sobivam valik piiratud eelarvega üliõpilase jaoks. Kombinatsioon paistis silma akadeemilise stiili, juhiste järgimise ja teksti struktureerimise poolest ning on oluliselt soodsam kui Opus. Pikendatud mõtlemisrežiim parandab keerulisemate peatükkide kvaliteeti.

\textbf{Lühikeste lõikude jaoks: GPT-5.4.} GPT-5.4 sobib hästi lühikeste tekstilõikude genereerimiseks ja keelelise korrektuuri jaoks. Kui soov on genereerida tervikdokumenti, siis ChatGPT vestlusliidese süsteemijuhised ja ohutuspoliitika võivad selle blokeerida või suunata ingliskeelsele väljundile, samas kui ChatGPT Codexi laiendus VS Code'is andis käesolevas katses sama viibaga tervikliku eestikeelse dokumendi (vt peatükk~\ref{subsec:gpt-toolo}).

\textbf{Kombineeritud praktika.} Iteratiivse juhtumiuuringu praktilise kogemuse põhjal tasub kaaluda tööjaotust mitme mudeli vahel: mahukamateks sisulisteks ja struktuurseteks ülesanneteks tugevaim kättesaadav arutlusmudel (käesolevas töös Opus 4.7 koos Claude Code'iga), keelekorrektuuriks ja kirjavigade otsimiseks aga odavam ja kiirem mudel (käesolevas töös GPT-5.5 Codexi laienduse kaudu).

\textbf{Autori maksimaalne kuutellimuse soovitus: ChatGPT Plus ja Claude Max 5\texttimes{} paralleelselt.} Üliõpilasele, kes soovib käesolevas töös kirjeldatud iteratiivset tööjaotust kahe pakkuja tugevamate mudelite vahel täismahus rakendada, soovitab autor pidada paralleelselt ChatGPT Plus (ca 23~€/kuu) ja Claude Max 5\texttimes{} (ca 106~€/kuu) tellimust ehk kokku ca 129~€/kuu (vt tabel~\ref{tab:mudelite-tellimused}). ChatGPT Plus annab GPT-5.4 ja Codex-laienduse piisava kasutuspiirangu juures keelekorrektuuri ja lühemate genereerimisülesannete jaoks, samas kui Claude Max 5\texttimes{} võimaldab laiemat Opus 4.6 ligipääsu ja Claude Code töövoo pikemaid sessioone ilma kvoodipiiranguid sageli ületamata.

\textbf{Autori minimaalne kuutellimuse soovitus: ChatGPT Plus ja Claude Pro.} Piiratud eelarvega üliõpilasele piisab käesoleva uurimuse kogemuse põhjal paralleelsest ChatGPT Plus (ca 23~€/kuu) ja Claude Pro (ca 22~€/kuu) tellimusest ehk kokku ca 45~€/kuu (vt tabel~\ref{tab:mudelite-tellimused}). Selle kooslusega saab teha kogu põhilise iteratiivse kirjutamis- ja toimetamistöö, kuigi pikemate sessioonide jooksul võivad mahupiirangud ette tulla.

\textbf{Kuutellimus on API-st oluliselt soodsam.} Käesoleva töö hinnakujunduse võrdlus (vt peatükk~\ref{subsec:hinnakujundus}) ja autori kasutuskogemus näitavad, et kuutellimusega töötamine on valdava enamuse üliõpilaste jaoks API kasutamisest mitmekordselt odavam. Kuutellimus annab fikseeritud ja prognoositava kuukulu (käesoleva töö soovitatud koosluste vahemikus ca 45--129~€/kuu), samas kui sama mahuga iteratiivse töövoo katmine API kaudu võib maksta sõltuvalt mudelist ja konteksti pikkusest mitusada eurot ning eeldab pidevat kuluseiret. API-päringuid tasub eelistada üksnes spetsiifilistel juhtudel, näiteks reprodutseeritavate empiiriliste hindamiste või automatiseeritud kvaliteedikontrolli puhul, kus iga päringu täpne kontekst ja kulu peavad olema dokumenteeritud.

\subsection{Soovitused lõputöö genereerimiseks agentsete TI tööriistadega} \label{subsec:agentsed-soovitused}

Käesoleva uurimuse praktiliste kogemuste põhjal, sealhulgas Opus 4.6 mudeliga Coworki kaudu genereeritud paroolihalduri lõputöö (vt peatükk~\ref{subsec:opus-toolo}) ning Opus 4.7 mudeliga Claude Code'i kaudu läbi viidud iteratiivne juhtumiuuring (vt peatükk~\ref{subsec:opus47-juhtum}), saab anda järgmised soovitused agentsete TI tööriistade kasutamiseks lõputöö genereerimisel. Siin peetakse silmas tööriistu, mis suudavad projekti faile lugeda ja muuta, käsurea käske käivitada ning \LaTeX{}-projekti struktuuri järgida. Eraldi tähelepanu väärib peatükis~\ref{subsec:opus47-juhtum} esitatud juhtum, mis näitab, et ühe tööpäeva pikkune juhitud iteratiivne seanss võib anda olulisi täiendusi ühe viiba lähenemisele -- näiteks tegeliku, kompileeritava lähtekoodi genereerimise. Samas on oluline rõhutada, et ka selle töövoo juures jäi empiiriliste väidete (näiteks raporteeritud F1-skoorid) sõltumatu kontroll endiselt inimese ülesandeks: tehniline reprodutseeritavus ei asenda sisulist auditeerimist.


\subsubsection{Genereerimise töövoog}

Soovitame järgmist samm-sammulist töövoogu:

\begin{enumerate}
    \item \textbf{Kasuta \LaTeX{}-malli algusest peale.} Paigalda ülikooli \LaTeX{}-mall (nt UniTartuCS) projektikausta enne genereerimise alustamist. Failipõhise ligipääsuga agentne tööriist saab seejärel malli struktuuri tuvastada ja järgida.
    \item \textbf{Genereeri sisukord ja struktuur.} Alusta sisukorra ja peatükkide ülesehituse genereerimisega. Anna mudelile ette teema, uurimisküsimused ja mahunõuded. Kontrolli ja kinnita struktuur enne edasi liikumist.
    \item \textbf{Genereeri peatükid järjest.} Genereeri iga peatükk eraldi, alustades sissejuhatusest. Iga järgmise peatüki genereerimisel on eelnevad peatükid kontekstis, mis tagab sidususe.
    \item \textbf{Genereeri tabelid ja joonised eraldi.} Tabelite ja jooniste \LaTeX{}-kood on keerulisem ja vajab eraldi tähelepanu. Genereeri need eraldi sammuna ja kontrolli visuaalset tulemust.
    \item \textbf{Lisa viited ise või kontrolli neid põhjalikult.} Ära lase mudelil viiteid ``välja mõelda''. Viidete hallutsineerimise riski vähendamiseks järgi seda malli:
    \begin{itemize}
        \item otsi ise allikaid;
        \item lisa leitud allikate andmed BibTeX-faili;
        \item palu mudelil viiteid tekstis kasutada ainult olemasolevatest allikatest.
    \end{itemize}
    \item \textbf{Kompileeri regulaarselt.} Kasuta tööriista või arenduskeskkonna võimalust käivitada \texttt{pdflatex} ja \texttt{biber} käske, et kontrollida, kas genereeritud \LaTeX{}-kood kompileerub vigadeta.
    \item \textbf{Vaata üle ja paranda.} Iga genereeritud peatüki järel vaata tekst üle, märgi probleemkohad ja palu mudelil need parandada. Ära jäta ülevaatamist töö lõppu.
\end{enumerate}

\subsubsection{Parimad praktikad}

\begin{itemize}
    \item \textbf{Kasuta võimaluse korral pikendatud mõtlemisrežiimi.} Pikendatud mõtlemisrežiim või muu sügavama arutluse seadistus (\textit{extended thinking}) võib parandada keeruliste akadeemiliste tekstide kvaliteeti. Sügavam arutlus on eriti kasulik sissejuhatuse, metoodika ja arutelu peatükkide jaoks.
    \item \textbf{Anna konkreetseid juhiseid.} Mida konkreetsemad on juhised, seda parem on tulemus. Üldise palve ``kirjuta sissejuhatus'' asemel täpsusta töö maht ja peamised teemad.
    \item \textbf{Kontrolli viiteid alati.} Ükski TI mudel ei genereeri 100\% korrektseid viiteid alati. Kontrolli iga viite olemasolu ja korrektsust käsitsi.
    \item \textbf{Ära usalda ühe kasutajaviibaga (\textit{single-prompt}) tulemust pimesi ilma kriitilise ülevaatamiseta.} Käesoleva uurimuse eksperimendis andis ühe viiba lähenemine mõnel juhul tugeva vormilise tulemuse, kuid sisaldas endiselt hallutsineeritud viiteid ja kontrollimata empiirilisi väiteid (vt peatükk~\ref{eksperimendid}). Märkimisväärselt parema tulemuse saavutamiseks on mõistlik kombineerida ühe viibaga genereerimist iteratiivsete paranduste ja kasutaja ülevaatustega.
    \item \textbf{Ole kursis enda tööga.} TI on abivahend, mitte autor. Üliõpilane peab iga genereeritud lõigu sisuliselt üle vaatama ja vajadusel ümber kirjutama.
    \item \textbf{Kasuta keeruliste peatükkide jaoks tugevamat mudelit.} Keerukamate peatükkide, näiteks metoodika ja arutelu, genereerimiseks või põhjalikuks toimetamiseks tasub valida hetkel kättesaadav tugevam arutlusmudel. Lihtsamate korrektuuri- ja vormistusülesannete jaoks piisab sageli odavamast või kiiremini vastavast mudelist. Kuna mudeliperekonnad muutuvad kiiresti, tasub soovituse rakendamisel kontrollida, millised keelemudelid on praegu kättesaadavad.
\end{itemize}

TI ei korva üliõpilase sisulist tööd allikate, metoodika ja järelduste osas, kuid võimaldab korrektuuri- ja vormistustööd vähema mehaanilise töö ja ajakuluga teha.

\subsection{Tööriistad, mida vältida} \label{subsec:valtida}

Lisaks põhieksperimendis osalenud nelja mudeli võrdlusele testis autor lühemate mitteformaalsete katsete raames ka mõnda teist tööriista, mida lõputöö kirjutamise abivahendina ei saa soovitada, vaatamata nende laialdasele kättesaadavusele üliõpilastele. Järgmised hinnangud ei põhine käesoleva uurimuse formaalsel võrdlusel, vaid autori praktilisel kasutuskogemusel ning nende arendajate avalikul dokumentatsioonil.

\textbf{Perplexity ei sobi mahukate kirjutamisülesannete põhitööriistaks.} Perplexity on eelkõige TI-põhine otsingu- ja vastamissüsteem, mis sobib hästi üksikküsimuste, allikate esmase leidmise ja lühikeste kokkuvõtete jaoks. Lõputöö tervikpeatükkide genereerimisel ilmnes aga piirang: Perplexity kasutuskogemus ei ole samaväärne alusmudeli otse kasutamisega. Perplexity ametlik API-dokumentatsioon eristab otsingukonteksti suurust mudeli kontekstiaknast: otsingukontekst määrab, kui palju veebist leitud infot päringusse kaasatakse, samas kui kontekstiaken määrab maksimaalse tokenite hulga, mida mudel ühe päringu käigus töödelda saab \cite{perplexity2026pricing}. Agent API toetab küll eri pakkujate mudeleid \cite{perplexity2026models}, kuid Perplexity otsingupõhine töövoog rakendab lisaks otsingu, tsiteerimise, kulude ja päringu töötlemise piiranguid. Autori veebiliidese katsetes ei olnud võimalik pikemat töömustandit sama järjepidevalt kontekstis hoida nagu mudelite natiivsetes või failipõhistes agentsetes keskkondades. Ka DataStudiosi ülevaade rõhutab, et Perplexity failide ja vestluste puhul ei lisata kogu materjali automaatselt vastuse koostamise hetkel mudeli konteksti, vaid süsteem valib ja lisab aktiivsesse päringusse ainult asjakohaseid lõike \cite{datastudios2025perplexity}. Sellisest valikulisest kontekstist tulenevalt võib pikema vestluse või suurema dokumendimaterjali korral kaduda osa varem kokkulepitud tingimustest, struktuurist või piirangutest. Kokkuvõttes sobib Perplexity hästi infootsinguks ja lühikeste küsimuste lahendamiseks, kuid mitte lõputöö genereerimise põhitööriistaks, kus on vaja mahukat, püsivat ja täpselt kontrollitavat projektikonteksti.

\textbf{GitHub Copilot Pro ei sobi akadeemilise teksti genereerimiseks.} GitHub Copilot Pro, mida Tartu Ülikool pakub üliõpilastele tasuta GitHub Education programmi kaudu, on mõeldud eelkõige koodi kirjutamise abivahendiks \cite{github2026copilot}. GitHubi enda dokumentatsiooni kohaselt on Copilot Chat mõeldud ainult kodeerimisega seotud päringutele ning juhiste piiramine kodeerimisülesannetega parandab mudeli väljundi kvaliteeti \cite{github2026copilot_responsible}. Lõputöö genereerimise kontekstis osutus kodeerimisele orienteeritud Copilot aga ebapiisavaks. Copiloti mudelid tundusid oluliselt piiratumad -- nad kippusid andma lühikesi, pealiskaudseid vastuseid ja vältisid pikemate akadeemiliste tekstilõikude genereerimist. Nguyen ja Nadi \cite{nguyen2022empirical} on oma empiirilises uurimuses näidanud, et Copiloti soovitustel on tähemärkide pikkuse piirang, mis takistab pikema väljundi genereerimist, ning keerulisemate ülesannete korral langeb soovituste kvaliteet märgatavalt. Lisaks rakendab GitHub Copilotile päringupõhiseid limiite (\textit{rate limits}), mis piiravad kasutamise intensiivsust \cite{github2026copilot_ratelimits}. Kokkuvõttes on tööriist optimeeritud lühikeste koodilõikude ja kooditäienduste jaoks, mitte terviklike akadeemiliste tekstide loomiseks. Eeltoodud piirangute tõttu ei saa GitHub Copilot Pro-d soovitada lõputöö kirjutamise abivahendina, vaatamata tasuta kättesaadavusele üliõpilastele.

\subsection{Tulevikuväljavaated} \label{subsec:tulevikuvaljavaated}

TI mudelite ja neid ümbritsevate tööriistade areng on kiire ning käesoleva töö tulemused peegeldavad 2025--2026. aasta seisu. Lähitulevikus on tõenäoliselt olulised järgmised arengusuunad:

\begin{itemize}
    \item \textbf{Eesti keele toe sihipärane parandamine} -- eesti keele kvaliteedi paranemine ei sõltu ainult sellest, kui suur osa üldmudelite treeningandmetest on eestikeelne. Olulisemaks võivad osutuda kvaliteetsemad eestikeelsed terminibaasid, jätkutreenimine ja väikeste keelte jaoks kohandatud hindamismetoodikad. EstLLM projekti \cite{dorkin2026estllm} tulemused näitavad, et spetsialiseeritud jätkutreenimine võib eesti keele kvaliteeti märkimisväärselt tõsta ka olemasolevates mudelites.
    \item \textbf{Agentsete akadeemiliste töövoogude levik} -- Claude Code'i, Coworki, Google'i Antigravity ja OpenAI ChatGPT Codexi laienduse laadsed tööriistad näitavad, et akadeemilise kirjutamise tugi liigub vestlusaknast üha enam projektipõhisesse keskkonda. Sellise töövoo väärtus ei seisne ainult teksti genereerimises, vaid ka failide haldamises, \LaTeX{}-koodi parandamises, kompileerimisvigade diagnoosimises ja viidete tehnilises kontrollis. Käesolev eksperiment näitas, et agentne töövoog parandas genereerimise kvaliteeti kõigi testitud mudelite puhul.
    \item \textbf{Konteksti haldamise paranemine} -- kontekstiaknad võivad jätkuvalt suureneda, kuid ainult suurem tokenite arv ei taga paremat akadeemilist tulemust. Sama oluline on see, kuidas tööriist valib, kokku võtab ja säilitab varasemaid töö osi, allikaid, juhendeid ja kasutaja parandusi. Seetõttu muutub tõenäoliselt olulisemaks projektimälu, otsingupõhise konteksti ja failipõhise töövoo kvaliteet.
    \item \textbf{Valdkonnaspetsiifilised akadeemilised abivahendid} -- lähiaastatel võib kasvada tööriistade hulk, mis on kohandatud konkreetsete erialade, viitamisstiilide, ülikoolide juhendite või keelte jaoks. Sellised süsteemid ei pea tingimata olema eraldi suured mudelid; sama eesmärki võivad täita ka üldmudelid koos kontrollitud allikabaasi, terminiloendi, mallide ja automaatsete kvaliteedikontrollidega.
    \item \textbf{TI kasutamise läbipaistvam integreerimine õppetöösse} -- akadeemilised institutsioonid liiguvad tõenäoliselt pelgalt keelamise ja tuvastamise loogikalt TI kasutuse läbipaistva kirjeldamise ning hindamise suunas. See nõuab hindamismetoodikaid, mis keskenduvad üliõpilase sisulisele panusele, allikakriitikale ja tööprotsessi dokumenteerimisele, mitte ainult lõpptulemuse vormilisele sarnasusele TI-genereeritud tekstiga.
\end{itemize}

Käesolev uurimus piirdus neljast mudelist ja kolmest agentsest tööriistast koosneva võrdlusega 2025--2026. aasta seisuga. Arvestades TI valdkonna arengutempot, kus uusi mudeleid avaldatakse mitu korda aastas, tuleks sarnaseid teostatavusuuringuid korrata regulaarselt. See võimaldaks hinnata, kuidas ülalnimetatud arengud tegelikult realiseeruvad ja millised käesoleva töö järeldustest jäävad kehtima ka uuema põlvkonna mudelite puhul.

